\documentclass[12pt,a4paper]{article}
\usepackage[UTF8]{ctex}
\usepackage{amsmath, amssymb}
\usepackage{geometry}
\usepackage{booktabs}
\usepackage{enumitem}
\usepackage{tcolorbox}
\usepackage{float}
\usepackage{hyperref}

\geometry{left=2.5cm, right=2.5cm, top=2.5cm, bottom=2.5cm}

\title{\textbf{问题四 模型建立与求解（论文手工作稿）}\\[6pt]
\large VLSI 布图规划设计 --- L/T 型模块的最小包围盒（穷举 + 最优性证明）}
\author{2026 年第七届"华数杯"大学生数学建模竞赛 B 题}
\date{2026/08}

\begin{document}
\maketitle
\tableofcontents
\newpage

% ============================================================
\section{问题重述与模型建立}
% ============================================================

\subsection{问题重述}

\subsubsection*{背景与意义}

在实际芯片设计中，模块并非全部为矩形：电源域、模拟宏单元、异形 IP 常呈
L 型、T 型等多边形形状。这类模块无法直接用"宽 × 高"描述，但可视为
\textbf{若干矩形按固定相对位置拼接}而成。问题 4 在问题 1 的矩形布图模型
基础上引入 L/T 型模块与四向旋转，检验"矩形分解 + B$^*$ 树"框架对
\textbf{非矩形模块}的推广能力，是模型修正与算法延续性的验证题。

\subsubsection*{题目重述}

问题 4 要求：给定 4 个模块——T 型（b1）、L 型（b2）与两个矩形
（b3 $2\times1$、b4 $1\times4$）——每个模块可进行
$0^\circ / 90^\circ / 180^\circ / 270^\circ$ 四种旋转，求无重叠布图使
\textbf{包围盒面积最小}，并给出最优摆放与旋转角。

\subsubsection*{模块几何（矩形分解）}

L/T 型模块表示为\textbf{固定相对位置的矩形并集} $M = \bigcup_{k} R_k$：
\begin{table}[H]
\centering
\caption{模块定义（矩形分解，总面积 24）}\label{tab:mod}
\small
\begin{tabular}{clcl}
\toprule
模块 & 形状 & 分解矩形（相对坐标） & 面积 \\
\midrule
b1 & T 型 & 顶横条 $4\times2$ @(0,2)；竖干 $2\times2$ @(1,0) & 12 \\
b2 & L 型 & 左竖条 $1\times4$ @(0,0)；右下横条 $1\times2$ @(1,0) & 6 \\
b3 & 矩形 & $2\times1$ & 2 \\
b4 & 矩形 & $1\times4$ & 4 \\
\bottomrule
\end{tabular}
\end{table}

例如 b1（T 型）为 $4\times2$ 顶横条与居中 $2\times2$ 竖干的并集：
外接 bbox 为 $4\times4$，但实际面积仅 12（存在 $4\times4$ 的凹角空缺）。
b2（L 型）外接 bbox 为 $2\times4$，实际面积 6。四个模块总面积
$\Sigma A = 24$。

\subsubsection*{数学形式化}

每个模块经四向旋转（$90^\circ$ 倍角，角点变换 $(x,y) \mapsto (-y,x)$
后归一化平移）后共有 4 种朝向；模块整体作为\textbf{超块}（内部矩形相对
位置固定）参与布局。设包围盒宽高为 $(W, H)$，问题形式化为
\begin{equation}
\min W \times H, \qquad \text{s.t. 全部超块互不重叠（矩形级）、无越界}.
\end{equation}

\textbf{面积下界}：任意合法布局的包围盒面积 $\ge$ 模块总面积，
\begin{equation}
WH \ge \Sigma A = 24,
\end{equation}
该下界可判优：若找到面积恰为 24 的合法布局，即为全局最优。

\subsubsection*{目标与问题性质}

目标：$\min WH$（包围盒面积）；约束：矩形级无重叠。本问与问题 1 的
关系是\textbf{模型修正而非重新建模}——仅表示粒度（矩形 → 超块）与
解码支撑（凹角精确支撑）发生改变，目标函数与搜索框架不变。由于实例
仅 4 模块，拓扑空间极小：14 树形 $\times$ 24 排列 $\times$ 256 旋转
$= 86\,016$ 组合，可\textbf{穷举全部组合}并配合面积下界严格\textbf{证明
全局最优}——结果确定，不依赖任何搜索运气（这与问题 1--3 的启发式
求解形成方法上的对照）。

\subsection{问题分析}

问题 4 是问题 1 模型的\textbf{模型修正 + 算法延续}验证：

\begin{enumerate}
    \item \textbf{两处变化}：模块不再全为矩形（出现 L/T 型）；
    旋转扩展为四向（$90^\circ$ 倍角）；
    \item \textbf{延续性}：非矩形模块经\textbf{矩形分解}后，B$^*$ 树表示
    框架与轮廓解码框架不变，仅节点语义与解码支撑扩展——满足题目
    "在问题 1 算法基础上适当修改，不要求重新设计新的优化算法"；
    \item \textbf{规模决定方法}：实例仅 4 模块，拓扑空间极小
    （14 树形 $\times$ 24 排列 $\times$ 256 旋转 $= 86\,016$ 组合），
    直接\textbf{穷举全部组合}并配合面积下界\textbf{证明全局最优}，
    不依赖搜索运气。
\end{enumerate}

\subsection{模型修正}

\subsubsection{修正一：模块表示 --- 矩形分解}

L/T 型模块可表示为\textbf{固定相对位置的矩形并集}：
\begin{equation}
M = \bigcup_{i=1}^{k} R_i, \qquad
R_i = (x_i, y_i, w_i, h_i) \text{ 相对坐标固定}.
\end{equation}

例如 T 型（b1）：顶横条 $4\times2$ 与居中竖干 $2\times2$
（bbox $4\times4$，面积 12）；L 型（b2）：左竖条 $1\times4$ 与
右下横条 $1\times2$（bbox $2\times4$，面积 6）。

\subsubsection{修正二：节点语义 --- 矩形升级为超块}

B$^*$ 树节点从"单个矩形"升级为"\textbf{超块}"（矩形组，内部相对位置
固定、整体移动与旋转）。树的父子/左右子关系语义\textbf{不变}。

\subsubsection{修正三：解码 --- 凹角精确支撑}

超块放置时，其 $y$ 坐标由内部各矩形分别对轮廓求支撑、取最大值：
\begin{equation}
\label{eq:support}
y_{\text{block}} = \max_{i} \big( \text{support}(R_i) - y_i \big),
\end{equation}
其中 $\text{support}(R_i)$ 为 $R_i$ 的 $x$ 区间与已放置矩形相交的
最高顶。该式使 \textbf{L/T 凹角空间被充分利用}（对比"按 bbox 放置"
浪费凹角），是无重叠前提下更紧凑的解码。

\subsubsection{修正四：旋转 --- 四向}

$0^\circ/90^\circ/180^\circ/270^\circ$ = 分解矩形组的整体旋转：对组内
所有矩形做角点变换 $(x, y) \mapsto (-y, x)$ 并归一化平移，\textbf{预计算
4 组坐标}供穷举使用。

\subsection{目标与约束}

与问题 1 完全一致：目标为\textbf{包围盒面积最小}，约束为\textbf{模块
（矩形级）互不重叠}。修正仅发生在\textbf{表示粒度}与\textbf{解码支撑}，
目标函数与搜索框架未变——这是"模型修正而非重新建模"的关键论证。

% ============================================================
\section{求解算法设计}
% ============================================================

\subsection{总体框架}

\begin{tcolorbox}[colback=blue!5, colframe=blue!60, title=算法流程]
\begin{enumerate}
    \item 预计算 4 个模块的四向旋转分解坐标（角点变换 + 归一化）；
    \item 穷举 B$^*$ 树拓扑：14 树形 $\times$ 24 标签排列；
    \item 每拓扑下穷举 $4^4 = 256$ 旋转组合；
    \item 超块解码（凹角精确支撑）$\rightarrow$ 矩形级重叠检查
    $\rightarrow$ 更新最小面积 best；
    \item 对照：按超块 bbox 放置的\textbf{必合法版本}求最小值。
\end{enumerate}
\end{tcolorbox}

\subsection{穷举空间}

\begin{equation}
14 \;(\text{树形}) \times 24 \;(\text{排列}) \times 256 \;(\text{旋转})
= 86\,016 \text{ 组合}.
\end{equation}

\begin{itemize}
    \item 树形：$C_4 = 14$ 种有序二叉树形状（递归生成全部相对索引组合）；
    \item 排列：4 模块分配到树节点，$4! = 24$；
    \item 旋转：每模块 4 向，$4^4 = 256$。
\end{itemize}

\subsection{解码与合法性检查}

\begin{itemize}
    \item 超块解码：$x$ 由树关系确定（左子紧贴父右、右子与父左对齐），
    $y$ 按式 \eqref{eq:support} 凹角精确支撑；
    \item \textbf{矩形级重叠检查}：全部 86\,016 布局逐对验证
    6 个分解矩形两两无重叠（不信任解码，独立检查）；
    \item 剪枝与断言：$best\_area \ge \Sigma A = 24$（下界恒成立）、
    $best\_area \le bbox\_version$（精确版不劣于 bbox 版）。
\end{itemize}

% ============================================================
\section{最优性证明与结果}
% ============================================================

\subsection{最优性证明}

\textbf{命题}：最小包围盒面积为 24，且为全局最优。

\textbf{证明}：
\begin{enumerate}
    \item \textbf{下界}：任意合法布局的包围盒面积 $\ge$ 模块总面积
    $= 24$（4 模块互不重叠，任何包围盒必须覆盖全部模块面积）；
    \item \textbf{可达性}：穷举全部 86\,016 组合，找到面积 $= 24$ 的
    合法布局（矩形级无重叠验证通过）；
    \item 由 1、2：$24 \le$ 最小面积 $\le 24$，故最小面积 $= 24$，
    \textbf{全局最优}（下界可达性证明）。
\end{enumerate}

该证明\textbf{不依赖搜索运气}：穷举覆盖全部 B$^*$ 树拓扑（完备性），
任何更优布局（面积 $< 24$）将违反下界，故不存在。

\subsection{最优摆放}

\begin{table}[H]
\centering
\caption{最优布局（包围盒 $4 \times 6 = 24$）}
\small
\begin{tabular}{lcccc}
\toprule
模块 & 旋转角 & 位置 $(x,y)$ & 分解矩形（绝对坐标） \\
\midrule
b3 & $0^\circ$ & (0,0) & $2\times1$ @(0,0) \\
b1 & $270^\circ$ & (0,0) & 竖干 $2\times4$ @(2,0)；横条 $2\times2$ @(0,1) \\
b2 & $270^\circ$ & (0,3) & 横条 $4\times1$ @(0,4)；横条 $2\times1$ @(0,3) \\
b4 & $90^\circ$ & (0,5) & $4\times1$ @(0,5) \\
\bottomrule
\end{tabular}
\end{table}

\begin{itemize}
    \item \textbf{最小包围盒}：$4\times6 = 24$（$=$ 模块总面积下界，
    全局最优，完美铺满无死区）；
    \item b3（$2\times1$）恰好填入 b1（T 型 $270^\circ$）的\textbf{凹角}，
    体现凹角精确支撑的解码优势；
    \item \textbf{对照}：按超块 bbox 放置的最小面积为 32（$>24$），
    证明精确支撑语义的有效性（省 25\%）；
    \item 穷举 86\,016 组合全部通过矩形级重叠检查（合法率 100\%）。
\end{itemize}

\subsection{结果讨论}

\begin{itemize}
    \item 达到理论下界（$= $ 模块总面积）意味着\textbf{零死区完美铺满}，
    这是矩形打包问题的理想情形，本实例恰好可达；
    \item 模型修正的有效性由 bbox 对照（32 vs 24）量化证明；
    \item 大规模非矩形实例可沿用问题 1 的模拟退火流程，仅将解码替换
    为超块版本（本问框架的直接推广）。
\end{itemize}

% ============================================================
\section{正确性保障与可复现性}
% ============================================================

\subsection{工程保障}

\begin{itemize}
    \item \textbf{矩形级重叠检查}：全部 86\,016 布局逐对验证
    6 矩形无重叠（不信任解码，独立检查）；
    \item \textbf{面积下界断言}：$best\_area \ge 24$ 恒成立；
    \item \textbf{bbox 版对照}：精确版（24）$\le$ bbox 版（32），
    防止支撑语义实现错误；
    \item \textbf{旋转正确性}：四向坐标由角点变换推导，可手工复核；
    \item 求解器为独立纯 Python 强算器（不依赖 C++ 核心），
    秒级完成，结果完全确定。
\end{itemize}

\subsection{运行命令与产物}

{\small
\begin{verbatim}
conda activate math
python q4/q4_solver.py                    # 强算 + 验证（秒级）
python q4/visualization/plot_q4.py \
    --csv q4/output/q4_result.csv \
    --out q4/visualization/figs/q4_optimal.png
\end{verbatim}
}

\begin{itemize}
    \item \texttt{q4/output/q4\_result.csv}：最优布局（min\_area 24、
    下界 24、bbox 版 32、合法布局 86016）；
    \item \texttt{q4/visualization/figs/q4\_optimal.png}：最优摆放示意图。
\end{itemize}

% ============================================================
\begin{thebibliography}{9}
\bibitem{bstar} Chang Y C, Chang Y W, Wu G M, et al. B*-Trees: a new
  representation for non-slicing floorplans[C]// Proceedings of the 37th
  Annual Design Automation Conference (DAC). 2000: 458--463.
\end{thebibliography}

\end{document}
